\section{The second-order theory}\label{sec:secondorder}

First order locates a crossing only up to a factor $1+o(1)$ in $T$. To locate
it up to $o(1)$ we need the constant in front of $N^{-(k-1)}e^{-T\lambda_F}$ for
every face that satisfies (H), together with the power of $T$. The plan is as
follows. Section~\ref{sec:tilt} defines the constant $K_F$.
Section~\ref{sec:locallaw} proves a local limit theorem inside a face and
deduces $M_F$ up to a factor $1+O(1/T+T^2/N)$. Section~\ref{sec:crossings}
compares 2 faces, Section~\ref{sec:papersAB} recovers Papers~A and~B, and
Section~\ref{sec:reducible} shows what can change without (H). In
this section gradients and Hessians in $\alpha$ are taken in the coordinates
$\alpha'=(\alpha_2,\dots,\alpha_k)$, $|\cdot|$ is the Euclidean norm, and
$q=(q_i)_{i\in F}$.

\subsection{The tilted chain and the constant}\label{sec:tilt}

To find the probability of a composition with fractions $\alpha$ we change the
rates so that $\alpha$ becomes typical. The optimal vector $r_\alpha$ in the
definition of $I_F(\alpha)$ gives a new generator $\widehat Q_\alpha$ on $F$, with
the same switching graph, whose stationary law is exactly $\alpha$. Under it the
visit fractions satisfy a central limit theorem with covariance
$\Gamma_\alpha/T$, and the price of the change of rates is $e^{-TI_F(\alpha)}$
times boundary factors for the start and the end of the path. The constant
$C_F(\alpha)$ collects these factors and the Gaussian normalization.

\begin{definition}[tilt and constants]\label{def:tilt}
Let $F$ be irreducible with $k\ge2$. For
$\alpha\in\Delta_F^\circ$ let $r_\alpha$ attain the supremum defining
$I_F(\alpha)$ (Proposition~\ref{prop:rate}(b)), and put
$\theta(\alpha)=A_Fr_\alpha/r_\alpha$ and $l_\alpha=\alpha/r_\alpha$
(entrywise). The \emph{tilted generator} is
$\widehat Q_\alpha=R^{-1}(A_F-\operatorname{diag}\theta(\alpha))R$ with
$R=\operatorname{diag}r_\alpha$. With $D=\operatorname{diag}\alpha$ let
\[
Z_\alpha=(\mathbf 1\alpha^\top-\widehat Q_\alpha)^{-1}-\mathbf 1\alpha^\top,
\qquad\Gamma_\alpha=DZ_\alpha+Z_\alpha^\top D,
\]
and let $\Sigma_\alpha$ be $\Gamma_\alpha$ with the row and column of state
$1$ deleted. Finally let $K_{\{i\}}=\mu_i$ for a single state, and
\[
C_F(\alpha)=\frac{(\mu\cdot r_\alpha)(l_\alpha\cdot\mathbf 1)}
{(2\pi)^{(k-1)/2}\sqrt{\det\Sigma_\alpha}},\qquad K_F=C_F(\alpha^*_F).
\]
\end{definition}

Nothing in Definition~\ref{def:tilt} changes when $r_\alpha$ is scaled. Note that
$l_\alpha\cdot r_\alpha=1$ and $I_F(\alpha)=\langle q-\theta(\alpha),\alpha\rangle$. Since
$r_\alpha$ minimizes $u\mapsto\sum_{i\ne j}\alpha_iq_{ij}u_j/u_i$, its derivative in $u_j$, which
is $\sum_il_{\alpha,i}q_{ij}-l_{\alpha,j}\theta_j(\alpha)$, vanishes at $r_\alpha$. Hence
\begin{equation}\label{eq:stationary}
l_\alpha^\top\bigl(A_F-\operatorname{diag}\theta(\alpha)\bigr)=0 .
\end{equation}
Also $(A_F-\operatorname{diag}\theta(\alpha))r_\alpha=0$ with $r_\alpha>0$, so $\theta(\alpha)\in
V_F$~\cite{bermanplemmons1994}. So $\widehat Q_\alpha$ is an irreducible generator with the
pattern of $A_F$ and, by~\eqref{eq:stationary}, stationary law $\alpha$. If
$(\mathbf 1\alpha^\top-\widehat Q_\alpha)v=0$, multiplying by $\alpha^\top$ gives $\alpha^\top
v=0$, so $\widehat Q_\alpha v=0$ and $v=0$. Thus $Z_\alpha$ exists, and
\begin{equation}\label{eq:Zprops}
Z_\alpha\mathbf 1=0,\qquad\alpha^\top Z_\alpha=0,\qquad
-\widehat Q_\alpha Z_\alpha=I-\mathbf 1\alpha^\top .
\end{equation}
So $\Gamma_\alpha\mathbf 1=0$. As $\Gamma_\alpha$ is symmetric, its adjugate is then a multiple of
$\mathbf 1\mathbf 1^\top$, so deleting another state instead of state $1$ gives the same
$\det\Sigma_\alpha$. At $\alpha^*_F$ the vector $u=r_F$ gives the value $\lambda_F=I_F(\alpha^*_F)$,
so $r_{\alpha^*_F}$ is a multiple of $r_F$, $\theta(\alpha^*_F)=q-\lambda_F\mathbf 1$, and
$l_{\alpha^*_F}$ is the matching multiple of $l_F$. In the local law below, $\mu\cdot r_\alpha$
accounts for the start of the path, $l_\alpha\cdot\mathbf 1$ for its free end, and $\Sigma_\alpha$
for the Gaussian spread of the occupation fractions.

For conductances $c_{ij}=c_{ji}$ on $F$ let
$\tree(c)=\sum_{\mathcal T}\prod_{\{i,j\}\in\mathcal T}c_{ij}$, the sum over
the spanning trees $\mathcal T$ of the complete graph on $F$.

\begin{proposition}[the constant]\label{prop:constant}
Let $F$ be irreducible with $k\ge2$.
\begin{enumerate}[(a)]
\item The maps $\alpha\mapsto\theta(\alpha)$, $I_F(\alpha)$, $\Sigma_\alpha$,
$r_\alpha$ and $l_\alpha$ (with $\sum_ir_{\alpha,i}=1$) are real-analytic on
$\Delta_F^\circ$, $\Sigma_\alpha$ is positive definite, and
$\nabla^2I_F(\alpha)=\Sigma_\alpha^{-1}$. Hence
$\det\Sigma_\alpha\cdot\det\nabla^2I_F(\alpha)=1$ and
$C_F(\alpha)=(\mu\cdot r_\alpha)(l_\alpha\cdot\mathbf 1)
\sqrt{\det\nabla^2I_F(\alpha)/(2\pi)^{k-1}}$.
\item Suppose that $\nu_iq_{ij}=\nu_jq_{ji}$ for $i,j\in F$ and some
$\nu\in(0,\infty)^F$. Then $r_\alpha$ is a multiple of $\sqrt{\alpha/\nu}$
and $l_\alpha$ of $\sqrt{\alpha\nu}$, and $\widehat Q_\alpha$ is reversible
with conductances $c_{ij}(\alpha)=\sqrt{\alpha_i\alpha_j}\,q_{ij}\sqrt{\nu_i/\nu_j}$.
Also $I_F(\alpha)=\sum_iq_i\alpha_i-\sum_{i\ne j}c_{ij}(\alpha)$, and
\[
\det\Sigma_\alpha=\frac{2^{k-1}(\prod_i\alpha_i)^2}{\tree(c(\alpha))},\qquad
C_F(\alpha)=\frac{(\mu\cdot r_\alpha)(l_\alpha\cdot\mathbf 1)}{(4\pi)^{(k-1)/2}}
\,\frac{\sqrt{\tree(c(\alpha))}}{\prod_i\alpha_i}.
\]
\end{enumerate}
\end{proposition}

\begin{proof}
The proof of (a) is in Appendix~\ref{app:torus} and uses Lemma~\ref{lem:perron}.

(b) Put $u=\sqrt{\alpha/\nu}$ and $l=\alpha/u=\sqrt{\alpha\nu}$. Then
$\sum_il_iq_{ij}=\sum_i\sqrt{\alpha_i\nu_i}\,q_{ij}$ and
$l_j(A_Fu)_j/u_j=\nu_j\sum_iq_{ji}\sqrt{\alpha_i/\nu_i}$ agree because $\nu_jq_{ji}=\nu_iq_{ij}$.
So $u$ is a critical point of $\sum_{i\ne j}\alpha_iq_{ij}u_j/u_i$, which is convex in $\log u$,
and $r_\alpha$ is a multiple of $u$. Then $\alpha_i(\widehat Q_\alpha)_{ij}=l_iq_{ij}u_j=
c_{ij}(\alpha)$, which is symmetric by reversibility, and $\alpha_i\theta_i(\alpha)=\sum_j
c_{ij}(\alpha)$ gives $I_F$.

For the determinant let $\mathcal L=-D\widehat Q_\alpha$, the Laplacian of the conductances. Fix
$\eta$ with $\eta_1=0$, and let $u=Z_\alpha\eta$ and $\zeta=D\eta-(\alpha\cdot\eta)\alpha$.
By~\eqref{eq:Zprops}, $\mathcal Lu=\zeta$ and $\alpha\cdot u=0$, so
$\eta^\top\Gamma_\alpha\eta=2(D\eta)^\top u=2\zeta^\top u$. Let $v=u-u_1\mathbf 1$. As
$\zeta\cdot\mathbf 1=0$ and $\mathcal L\mathbf 1=0$, we get $\zeta^\top u=\zeta^\top v$ and
$\mathcal Lv=\zeta$, and since $v_1=0$ the rows $2,\dots,k$ give $\mathcal L_{(1)}v'=\zeta'$. Here
$\mathcal L_{(1)}$ is $\mathcal L$ without row and column $1$, primes denote the coordinates
$2,\dots,k$, and $\zeta'=(D'-\alpha'\alpha'^\top)\eta'$ with $D'=\operatorname{diag}\alpha'$.
Hence $\Sigma_\alpha=2(D'-\alpha'\alpha'^\top)\mathcal L_{(1)}^{-1}(D'-\alpha'\alpha'^\top)$. The
matrix determinant lemma gives $\det(D'-\alpha'\alpha'^\top)=\prod_{i\ge2}\alpha_i\,
(1-\sum_{i\ge2}\alpha_i)=\prod_i\alpha_i$, and the matrix-tree theorem~\cite{bollobas1998} gives
$\det\mathcal L_{(1)}=\tree(c(\alpha))$.
\end{proof}

\begin{corollary}[symmetric faces]\label{cor:symface}
Let $F$ satisfy \textup{(H)} with $k\ge2$, and let $Q_F$ be symmetric with constant row sums.
Then $\lambda_F=q_i-\sum_{j\in F\setminus\{i\}}q_{ij}$ for every $i\in F$, $\alpha^*_F$ is uniform
on $F$, and
\[
K_F=\mu(F)\,k^{(k+1)/2}\sqrt{\tree(q_F)}\,(4\pi)^{-(k-1)/2},
\]
where $\tree(q_F)$ is $\tree(c)$ for the conductances $c_{ij}=q_{ij}$, $i,j\in F$.
\end{corollary}

\begin{proof}
The positive vector $\mathbf 1$ is a right and left eigenvector of $Q_F$, so it belongs to the
Perron root. So $\lambda_F$ is minus the row sum, $r_F$ and $l_F$ are multiples of $\mathbf 1$, and
$\alpha^*_F=\mathbf 1/k$.
At $\alpha=\mathbf 1/k$, Proposition~\ref{prop:constant}(b) with $\nu=\mathbf 1$ gives
$r_\alpha=\mathbf 1$, $l_\alpha=\alpha$ and $(\mu\cdot r_\alpha)(l_\alpha\cdot\mathbf 1)=\mu(F)$.
Also $c_{ij}(\alpha)=q_{ij}/k$, so $\tree(c(\alpha))=k^{-(k-1)}\tree(q_F)$, and
$\prod_i\alpha_i=k^{-k}$.
\end{proof}

\subsection{The local law and the face maxima}\label{sec:locallaw}

Inside the face the occupation fractions fluctuate on the scale $T^{-1/2}$ in $k-1$ directions. So
a single composition gets the factor $T^{(k-1)/2}$ times the volume $N^{-(k-1)}$ of a lattice cell.

\begin{theorem}[local law]\label{thm:local}
Let $F$ satisfy (H) with $k\ge2$, and let $\mathcal K\subset\Delta_F^\circ$ be
compact. Uniformly for $\alpha\in\mathcal K$ and $T\ge1$,
\[
f_T(\alpha)=T^{(k-1)/2}C_F(\alpha)e^{-TI_F(\alpha)}\bigl(1+O_{\mathcal K}(1/T)\bigr).
\]
For $F=\{i\}$, $f_T=\mu_ie^{-q_iT}$.
\end{theorem}

In words, the density of the visit fractions on the event that the chain stays
in $F$ is a Gaussian factor $T^{(k-1)/2}$ times the large-deviation factor
$e^{-TI_F(\alpha)}$ times the constant $C_F(\alpha)$, with a relative error
$O(1/T)$ that is uniform away from the boundary of the face.

\emph{Idea of the proof.} The proof is in Appendix~\ref{app:torus}. The sum
over words that defines $H$ is a matrix resolvent, so Cauchy's formula writes $H$
as an integral over a torus of radii $\theta(\alpha)$ (Lemma~\ref{lem:torus}).
All the coefficients are nonnegative, so the integrand is largest at the real
point, and a Laplace lemma with a parameter (Lemma~\ref{lem:laplace}) gives the
expansion. The Hessian at the real point is computed from the Perron root of a
perturbed generator (Lemma~\ref{lem:perron}), and this is where $\Sigma_\alpha$
comes from. The hypothesis $\mu(F)>0$ is needed, since $f_T=C_F=0$ when
$\mu(F)=0$. Also $f_T(\alpha)$ is continuous and positive in $(T,\alpha)$, as
$H$ is a convergent power series with some $W(m)>0$ by (H).

\begin{lemma}[discrete and continuum]\label{lem:discrete}
Let $F$ satisfy (H) with $k\ge2$, and let $\mathcal K\subset\Delta_F^\circ$ be
compact. Uniformly for $\operatorname{supp}n=F$ with $n/N\in\mathcal K$ and
$1\le T\le N^{1/2}$,
$P_N(n)=N^{-(k-1)}f_T(n/N)(1+O_{\mathcal K}(T^2/N))$.
\end{lemma}

The proof is in Appendix~\ref{app:torus}. With Theorem~\ref{thm:local} it sharpens the error in
Theorem~\ref{thm:first}(i) to $\frac{k-1}2\log T+\log C_F(n/N)+O(1/T+T^2/N)$.

\begin{theorem}[face maxima]\label{thm:facemax}
Let $F$ satisfy (H). If $k\ge2$, then uniformly for $1\le T\le N^{1/3}$,
\[
M_F=N^{-(k-1)}T^{(k-1)/2}K_Fe^{-T\lambda_F}\bigl(1+O(1/T+T^2/N)\bigr),
\]
and every maximizer $\hat n$ satisfies $\hat n/N=\alpha^*_F+O(1/T)$. For $k=1$,
$M_{\{i\}}$ is given by~\eqref{eq:vertexmax}. In both cases,
uniformly for $1\le T\le N^{1/3}$,
\begin{equation}\label{eq:logmax}
\log M_F=-(k-1)L+\tfrac{k-1}2\log T+\log K_F-T\lambda_F+O(1/T+T^2/N).
\end{equation}
\end{theorem}

In words, the best composition on $F$ sits within $O(N/T)$ of $N\alpha^*_F$,
and its probability is the Gaussian peak of the local law at $\alpha^*_F$. The
factor $N^{-(k-1)}$ is the volume of a lattice cell, and $T^{(k-1)/2}$ is the
height of a Gaussian density whose spread in each free direction is of order
$T^{-1/2}$.

\emph{Idea of the proof.} Near $\alpha^*_F$ we use the local law and
Lemma~\ref{lem:discrete}, and maximize a Gaussian with a small linear
perturbation. Away from $\alpha^*_F$ the rate exceeds $\lambda_F$ by a fixed
amount, and Lemma~\ref{lem:upper}, applied with finitely many vectors $\theta$,
shows that these compositions are exponentially less likely in $T$. For bounded
$T$ positivity and continuity suffice.

\begin{proof}
For $k=1$ this is~\eqref{eq:vertexmax}. Let $k\ge2$, put
$X=N^{-(k-1)}T^{(k-1)/2}K_Fe^{-T\lambda_F}$ and $x=\alpha'-\alpha^{*\prime}_F$. By
Propositions~\ref{prop:rate}(a) and~\ref{prop:constant}(a), $I_F$ and $\log C_F$ are
real-analytic near $\alpha^*_F$, $I_F$ has its minimum $\lambda_F$ there and $\nabla^2I_F>0$. So
there is a closed ball $U=\{|x|\le\rho_0\}\subset\Delta_F^\circ$ on which
$c_U|x|^2\le I_F-\lambda_F\le C_U|x|^2$, $|\log(C_F/K_F)-\gamma\cdot x|\le C_1|x|^2$ with
$\gamma=\nabla\log C_F(\alpha^*_F)$, and $\log C_F$ is $C_3$-Lipschitz.

\emph{Outside $U$.} As $I_F$ is lower semicontinuous, Proposition~\ref{prop:rate}(a) gives
$I_F\ge\lambda_F+3\eta$ on $\Delta_F\setminus U^\circ$ for some $\eta>0$. Since
$I_F(\beta)=\sup_{u>0}\langle q-A_Fu/u,\beta\rangle$ and $A_Fu/u\in V_F$, compactness gives
finitely many $\theta_j\in V_F$ with $\max_j\langle q-\theta_j,\beta\rangle>\lambda_F+2\eta$ on
$\Delta_F\setminus U^\circ$.
Lemma~\ref{lem:upper} then gives
$P_N(n)\le C(T/N)^{k-1}e^{-T(\lambda_F+2\eta)}=X\cdot O(T^{(k-1)/2}e^{-2\eta T})$ for
$\supp n=F$ and $n/N\notin U^\circ$.

\emph{Inside $U$.} By Theorem~\ref{thm:local} and Lemma~\ref{lem:discrete} with $\mathcal K=U$,
for $\alpha=n/N\in U$,
\begin{equation}\label{eq:insideU}
P_N(n)=X\,\tfrac{C_F(\alpha)}{K_F}\,e^{-T(I_F(\alpha)-\lambda_F)}(1+E),
\qquad|E|\le C_2(1/T+T^2/N).
\end{equation}
For $T\ge2C_1/c_U$ the middle factors are at most
$\exp(\sup_x[\gamma\cdot x-\frac{c_UT}2|x|^2])=e^{|\gamma|^2/(2c_UT)}$, so
$M_F\le X(1+O(1/T+T^2/N))$. For the lower bound let $n^0$ have support $F$ and
$|n^0_i-N\alpha^*_{F,i}|<1$ for all $i$, which exists for $N\ge N_0$ since
every $\alpha^*_{F,i}>0$. Then $C_F(n^0/N)\ge K_Fe^{-kC_3/N}$ and
$T(I_F(n^0/N)-\lambda_F)\le C_Uk^2T/N^2$, so~\eqref{eq:insideU} gives
$P_N(n^0)\ge X(1-O(1/T+T^2/N))$. Fix $T_2$ so large that the 2 bounds above give
$M_F/X\in[\frac12,\frac32]$ for $T\ge T_2$. For $1\le T\le T_2$, $f_T$ is positive and continuous
on $[1,T_2]\times U$, so Lemma~\ref{lem:discrete} gives $M_F/X\in[c,C]$. This gives the first
display and~\eqref{eq:logmax}.

\emph{The maximizer.} For $T\ge T_3$ with $T_3$ large, the bound outside $U$ is below $X/2\le
P_N(n^0)$, so $\hat\alpha=\hat n/N\in U$ and $|E|\le\frac12$. Taking logarithms in $P_N(\hat
n)\ge P_N(n^0)$ gives $c_UT|\hat x|^2\le C_3|\hat x|+C_4(1/T+T^2/N)\le C_3|\hat x|+2C_4/T$ with
$\hat x=\hat\alpha'-\alpha^{*\prime}_F$, because $T\le N^{1/3}$ gives $T^2/N\le1/T$. This step uses
the exponent $\frac13$, and it gives $|\hat x|=O(1/T)$. For $1\le T<T_3$ the claim holds since
$|\hat x|$ is bounded.
\end{proof}

\subsection{Crossings and tie windows}\label{sec:crossings}

We now compare 2 faces. Taking logarithms in Theorem~\ref{thm:facemax}, the
difference $\log M_F-\log M_G$ is an explicit increasing function of $T$ up to
$O(1/T+T^2/N)$, and its zero is the crossing.

\begin{theorem}[the crossing equation]\label{thm:crossing}
Let $F$ and $G$ satisfy (H), with $\Delta=k_F-k_G\ge1$ and
$\delta=\lambda_G-\lambda_F>0$, and let $D_N(T)=\log M_F-\log M_G$.
\begin{enumerate}[(i)]
\item For large $N$, $D_N$ has a zero in
$[\frac{\Delta}{2\delta}L,\frac{2\Delta}{\delta}L]$.
\item For large $N$, every zero $T_N\in[1,N^{1/3}]$ of $D_N$ satisfies
\begin{equation}\label{eq:crossing}
\delta T_N+\frac\Delta2\log T_N=\Delta L-\log\frac{K_F}{K_G}
+O\Bigl(\frac1L\Bigr).
\end{equation}
\item All such zeros lie in one interval of length $O(1/L)$, and
\[
T_N=\frac\Delta\delta\Bigl[L-\frac12\log L+b_{FG}\Bigr]
+O\Bigl(\frac{\log L}L\Bigr),\qquad
b_{FG}=-\frac12\log\frac\Delta\delta-\frac1\Delta\log\frac{K_F}{K_G}.
\]
\end{enumerate}
\end{theorem}

\begin{proof}
By~\eqref{eq:logmax}, uniformly on $[1,N^{1/3}]$,
$D_N(T)=-\Delta L+\frac\Delta2\log T+\delta T+\log\frac{K_F}{K_G}+O(1/T+T^2/N)$.

(i) At $T=\frac{\Delta}{2\delta}L$ this is $-\frac\Delta2L+O(\log L)<0$, and at
$T=\frac{2\Delta}\delta L$ it is $\Delta L+O(\log L)>0$. For fixed $N$, $M_F$ is a maximum of
finitely many polynomials in $h$, and it is positive because $F$ is admissible. So $D_N$ is
continuous and has a zero in between.

(ii) Let $\psi(T)=\delta T+\frac\Delta2\log T$, which is increasing. At a zero
$T_N\in[1,N^{1/3}]$, $\psi(T_N)=\Delta L-\log(K_F/K_G)+O(1)$. This is larger than
$\psi(\frac\Delta{2\delta}L)=\frac\Delta2L+O(\log L)$, and $\delta T_N\le\psi(T_N)$. So
$\frac\Delta{2\delta}L<T_N\le\frac{2\Delta}\delta L$, and then $O(1/T_N+T_N^2/N)=O(1/L)$.

(iii) Let $\psi(T^\circ)=\Delta L-\log(K_F/K_G)$. As $\psi'\ge\delta$, every zero is within
$O(1/L)$ of $T^\circ$. Also $T^\circ=\frac\Delta\delta L+O(\log L)$, so
$\log T^\circ=\log L+\log\frac\Delta\delta+O(\log L/L)$, which gives the explicit form.
\end{proof}

In words: the larger face $F$ pays $N^{-\Delta}$ for spreading and gains
$e^{\delta T}$ for staying, and its Gaussian peak gives an extra factor
$T^{\Delta/2}$; the crossing balances the 3. The explicit form in (iii)
converges slowly, so computations should solve~\eqref{eq:crossing} (see
Corollary~\ref{cor:C4}). We do not prove that the zero of $D_N$ is unique in
general.

Sometimes several faces tie at first order: at some $\tau^*$ more than 2 faces
minimize $\Phi_{\tau^*}$, as all faces do for the model of Paper~B. Then the
competition plays out in a window $T=\tau^*(L-\frac12\log L+t)$ with $t$ of
order $1$. In this window each tied face has a line $\ell_F(t)$, and the mode
follows the highest line.

\begin{theorem}[tie windows]\label{thm:window}
Fix $\tau^*>0$. Let $m^*$ be the least value of $\Phi_{\tau^*}$ over the
admissible faces and $\mathcal F^*$ the set of minimizers. Assume that every
face in $\mathcal F^*$ is irreducible. By Theorem~\ref{thm:hull}(c) this holds
if $Q$ has symmetric support or $\mu$ has full support. Put $T=\tau^*(L-\frac12\log L+t)$, and for
$F\in\mathcal F^*$ let
$\ell_F(t)=\log K_F+\frac{k_F-1}2\log\tau^*-\tau^*\lambda_Ft$ and
$E(t)=\max_{G\in\mathcal F^*}\ell_G(t)$.
\begin{enumerate}[(i)]
\item Uniformly for $t$ in compact sets, for $F\in\mathcal F^*$,
$\log M_F=-m^*L+\frac{m^*}2\log L+\ell_F(t)+O(\log L/L)$.
\item There are $\varepsilon_1,\varepsilon_0>0$ such that for large $N$ and
$|T/L-\tau^*|\le\varepsilon_1$, every admissible $F\notin\mathcal F^*$ has
$M_F\le N^{-\varepsilon_0}\max_{G\in\mathcal F^*}M_G$.
\item There are $\varepsilon,t_0>0$ such that for large $N$ and
$t_0\le|t|\le\varepsilon L$, every global mode lies on a face
$F\in\mathcal F^*$ with $\ell_F=E$ on $[t_0,\infty)$ if $t>0$, or on
$(-\infty,-t_0]$ if $t<0$. Such a face has the least $\lambda_F$ in
$\mathcal F^*$ if $t>0$ and the largest if $t<0$.
\end{enumerate}
Consequently, for every $\eta>0$ and large $N$, if $|t|\le\varepsilon L$ is at
distance at least $\eta$ from every kink of $E$, then every global mode lies
on a face whose line equals $E$ near $t$. So the line of the face that
carries the global mode changes only at
$T=\tau^*(L-\frac12\log L+t_j)+o(1)$, where $t_j$ runs over the kinks of $E$,
and it changes near each kink. Faces with identical lines (the same $k_F$,
$\lambda_F$ and $K_F$) tie at this order and give no crossing.
\end{theorem}

In words: inside the window the logarithm of every tied face maximum is the
same leading term plus the line $\ell_F(t)$, faces that do not tie are smaller by
a power of $N$, and far out in the window the face with the least (for $t>0$) or
largest (for $t<0$) exit rate wins.

\begin{proof}
Every $F\in\mathcal F^*$ is admissible and irreducible, so it satisfies (H) by
Lemma~\ref{lem:admissible}(b). Write $T=\tau^*L(1+y)$ with $y=(t-\frac12\log L)/L$, so that
$|y|\le\frac12$ for $|t|\le L/4$ and large $N$. By~\eqref{eq:logmax} and
$\tau^*\lambda_F+k_F-1=m^*$, uniformly for $|t|\le L/4$,
\begin{equation}\label{eq:windowline}
\log M_F=-m^*L+\tfrac{m^*}2\log L+\ell_F(t)+\tfrac{k_F-1}2\log(1+y)+O(1/L)
\qquad(F\in\mathcal F^*).
\end{equation}
For $t$ in a compact set, $\log(1+y)=O(\log L/L)$, which gives (i).

(ii) Let $g>0$ be the least value of $\Phi_{\tau^*}(F)-m^*$ over the admissible
$F\notin\mathcal F^*$. As $0\le\lambda_F\le q_{\max}$, for
$|\tau-\tau^*|\le\varepsilon_1=g/(4q_{\max}+4)$ we have $\Phi_\tau(F)-\Phi_\tau(G)\ge g/2$ for
such $F$ and all $G\in\mathcal F^*$. With $\tau=T/L$, Theorem~\ref{thm:first}(ii) gives
$\log M_F-\log M_G\le-\frac g2L+O(\log L)\le-\frac g4L$.

(iii) Take $t_0\ge1$ larger than every $|t_j|$ and $\varepsilon\le\min(\frac14,\varepsilon_1/(2\tau^*))$,
so that (ii) applies. Let $t\ge t_0$, let $\ell_F=E$ on $[t_0,\infty)$, and let $G\in\mathcal F^*$
with $\ell_G\ne\ell_F$, so that $\lambda_F\le\lambda_G$. As both faces are tied,
$\Delta=k_F-k_G=\tau^*(\lambda_G-\lambda_F)\ge0$. If $\Delta=0$, \eqref{eq:windowline} gives
$\log(M_F/M_G)=\log(K_F/K_G)+O(1/L)$ with $K_F>K_G$. If $\Delta>0$,
$\log(M_F/M_G)=\Delta t+\ell_F(0)-\ell_G(0)+\frac\Delta2\log(1+y)+O(1/L)$, where
$|\log(1+y)|\le2|y|\le t/2$ for large $N$, so it is positive if $t_0$ is large. The case $t\le-t_0$ is the same.

For the consequence, fix $\eta>0$ and let $S_\eta=[-t_0,t_0]\setminus\bigcup_j(t_j-\eta,t_j+\eta)$.
For $G\in\mathcal F^*$, $E-\ell_G$ is convex, piecewise linear and nonnegative with kinks only at
kinks of $E$. So its zero set is empty, a single kink of $E$, or an interval whose finite endpoints
are kinks of $E$, and on each component of $S_\eta$ either $\ell_G=E$ or $E-\ell_G\ge g_\eta>0$.
Then (i) and (ii) give the claim on $S_\eta$, and (iii) for $t_0\le|t|\le\varepsilon L$.
\end{proof}

\subsection{Papers A and B}\label{sec:papersAB}

We use the polynomials $S_{a,b}$ and $S_N$, the root $u_N$, the crossing $p_N$,
$z_N=N(1-p_N)$ and $c_0=\frac12\log(\pi/8)$ of Section~\ref{sec:notation}. The
first result shows that a single state against a symmetric pair is Kagey's
board at every $N$.

\begin{proposition}[an exact pair identity]\label{prop:pair}
Let $i\ne j$ with $q_{ij}=q_{ji}=a$, $q_i=q_j=q$ and $\mu_i=\mu_j>0$, and let
$F=\{i,j\}$. For every $N\ge2$, $1\le m\le N-1$ and $h>0$ such that $I+hQ$ is
stochastic and $qh<1$,
\[
\frac{P_N(me_i+(N-m)e_j)}{P_N(Ne_i)}=S_{m,N-m}(u),\qquad u=\frac{ah}{1-qh}.
\]
So the balanced compositions, with $m=\lfloor N/2\rfloor$ or
$m=\lceil N/2\rceil$, are the most likely ones with support $F$,
$M_F=M_{\{i\}}S_N(u)$, and $M_F>M_{\{i\}}$ exactly when $u>u_N$.
\end{proposition}

\begin{proof}
A path with this composition is a word in $i,j$ with $m$ letters $i$. Each stay has probability
$1-qh$ and each switch $ah$, so a word with $s$ switches has probability
$\mu_i(1-qh)^{N-1-s}(ah)^s=\mu_i(1-qh)^{N-1}u^s$, while $P_N(Ne_i)=\mu_i(1-qh)^{N-1}$. Sum over
the words. The bin polynomials of~\cite{paperA} grow coefficientwise as the bin moves toward the
middle, and $S_{m,N-m}=S_{N-m,m}$. So the balanced $m$ give the largest $S_{m,N-m}(u)$. Finally
$S_N$ increases strictly and equals $1$ only at $u_N$ (Section~\ref{sec:notation}).
\end{proof}

The identity in Proposition~\ref{prop:pair} is checked in Lean.

\begin{corollary}[Paper A]\label{cor:paperA}
Let $d=2$, $q_{12}=q_{21}=1$ and $\mu=(\frac12,\frac12)$, the walk of~\cite{paperA} with
$T=N(1-p)$. Then $K_{\{1\}}=K_{\{2\}}=\frac12$ and $K_{\{1,2\}}=\sqrt{2/\pi}$, so
$b_{FG}=c_0$ for $F=\{1,2\}$ and $G=\{1\}$. For every $N\ge2$, $M_{\{1,2\}}=M_{\{1\}}$ holds
exactly at $T=z_N$, and $z_N+\frac12\log z_N=L+c_0+O(1/L)$.
\end{corollary}

\begin{proof}
Corollary~\ref{cor:symface} with $k=2$, $\tree(q_F)=1$ and $\mu(F)=1$ gives $K_{\{1,2\}}$, and
then $b_{FG}=-\log(K_{\{1,2\}}/K_{\{1\}})=c_0$. Proposition~\ref{prop:pair} with $a=q=1$ gives
$M_{\{1,2\}}=M_{\{1\}}S_N(h/(1-h))$, and $h/(1-h)=u_N$ exactly when $p=p_N$. So for every $N$
the only zero of $D_N$ is $z_N$, and by Theorem~\ref{thm:crossing}(i) it lies in
$[L/2,2L]\subset[1,N^{1/3}]$ for large $N$. Then Theorem~\ref{thm:crossing}(ii) with
$\Delta=\delta=1$ gives the equation.
\end{proof}

This is the crossing equation of~\cite{paperA} without its terms $1/(8z_N)$ and $R_N$. The term
$1/(8z_N)$ is $(c_{\{1\}}-c_{\{1,2\}})/z_N$ in the notation of Remark~\ref{rem:nextterm}, whose
proof we omit, with $c_{\{1\}}=0$ and $c_{\{1,2\}}=-\frac18$.

\begin{corollary}[Paper B]\label{cor:paperB}
Let $Q=J-dI$, where $J$ is the all-ones matrix, and let $\mu$ be uniform, the model
of~\cite{paperB}. A face of $k$ states has $\lambda_F=d-k$, $\alpha^*_F$ uniform and
$K_F=\frac1dk^{k+1/2}(4\pi)^{-(k-1)/2}$. So $\Phi_1(F)=d-1$ for every face, and all faces tie at
$\tau=1$. The crossing of a face of $k\ge2$ states with a vertex has $\Delta=\delta=k-1$ and
$b_{FG}=a_k=\frac12\log(4\pi)-\frac{k+1/2}{k-1}\log k$, the constant of the vertex crossing
of~\cite{paperB}.
\end{corollary}

\begin{proof}
Each $Q_F$ is symmetric with row sums $-(d-k)$, and the complete graph on $F$ has $k^{k-2}$
spanning trees by Cayley's formula. So Corollary~\ref{cor:symface} with $\mu(F)=k/d$ gives
$\lambda_F$, $\alpha^*_F$ and $K_F$. A vertex has $K=\frac1d$ and $\lambda=d-1$, so
$b_{FG}=-\frac1{k-1}\log(dK_F)=a_k$. The arithmetic that gives $b_{FG}$ from $K_F$ here and in
Corollary~\ref{cor:paperA} is checked in Lean. At $\tau^*=1$, Theorem~\ref{thm:window} describes
the tie window, and~\cite{paperB} resolves this window in full.
\end{proof}

\subsection{Reducible faces}\label{sec:reducible}

\begin{remark}[scope of the $\frac12$]\label{rem:half}
The coefficient $\frac12$ of $\log\log N$ in Theorems~\ref{thm:crossing} and~\ref{thm:window} is
due to the power $T^{(k-1)/2}$ in Theorem~\ref{thm:facemax}, so it needs only (H). At a kink of
$E$ where the line of $G$ gives way to that of $F$ as $t$ increases, $\Delta=\tau^*\delta$ and
$t_j=b_{FG}$. So every change in a tie window is a crossing of Theorem~\ref{thm:crossing}. The
best composition of a reducible face can visit a state only once,
with no Gaussian window in that direction. Then the $\frac12$ can fail, and in the next example it
is $1$ and then $0$.
\end{remark}

\begin{proposition}[a directed $3$-cycle]\label{prop:threecycle}
Let $d=3$, $q_{12}=b>2$, $q_{23}=q_{31}=1$, all other rates $0$, and
$\mu_1=1$.
\begin{enumerate}[(i)]
\item The admissible faces are $\{1\}$, $\{1,2\}$ (reducible, $\lambda=1$) and
$\{1,2,3\}$ (irreducible, $\lambda=0$). Exactly
$M_{\{1\}}=(1-bh)^{N-1}$ and $M_{\{1,2\}}=bh(1-h)^{N-2}$, and the latter is
attained only at $n=(1,N-1,0)$. For the full face,
$\alpha^*=(1,b,b)/(1+2b)$, $\det\nabla^2I=(1+2b)^4/(3b^2)$ and
$K_{\{1,2,3\}}=(1+2b)^2/(2\sqrt3\,\pi b)$.
\item $M_{\{1\}}=M_{\{1,2\}}$ has exactly one root $T^{(1)}_N$ with $bh<1$. It
satisfies $(b-1)T+\log T=L-\log b+O(T^2/N)$, so
$T^{(1)}_N=\frac1{b-1}[L-\log L+\log\frac{b-1}b]+O(\log L/L)$, and the
coefficient of $\log\log N$ in the bracket is $1$.
\item Every root $T\in[L/2,2L]$ of $M_{\{1,2\}}=M_{\{1,2,3\}}$ satisfies
$T=L+\log\frac{2\sqrt3\,\pi b^2}{(1+2b)^2}+O(1/L)$, so the coefficient of
$\log\log N$ is $0$. These roots lie in an interval of length $O(1/L)$.
\item Both are changes of the global mode. For large $N$ and $T\le N^{1/3}$,
the mode is $(N,0,0)$ for $T<T^{(1)}_N$, it is $(1,N-1,0)$ between
$T^{(1)}_N$ and the smallest root in (iii), and it lies on the full face above
the largest root in (iii).
\end{enumerate}
\end{proposition}

In words: the face $\{1,2\}$ is reducible, and its best composition
$(1,N-1,0)$ visits state $1$ once and then stays in state $2$. Its maximum is
about $b(T/N)e^{-T}$, with the power $T^1$ where an irreducible pair would have
$T^{1/2}$. This changes the coefficient of $\log\log N$ at both crossings.

\begin{proof}
(i) The chain moves only along $1\to2\to3\to1$ from $1$, which gives the supports and the
components $\{1\}$, $\{2\}$ of $\{1,2\}$. A path with support
$\{1,2\}$ is $1^{n_1}2^{n_2}$, of probability $(1-bh)^{n_1-1}bh(1-h)^{n_2-1}$, and raising $n_1$
by one multiplies it by $(1-bh)/(1-h)<1$. For the full face $\lambda=0$, $r=\mathbf 1$, and
$l=\alpha^*=(1,b,b)/(1+2b)$ is the stationary law, so $(\mu\cdot r)(l\cdot\mathbf 1)=1$. The
ratios $u_2/u_1$, $u_3/u_2$ and $u_1/u_3$ have product $1$, so the AM--GM inequality gives
$I(\alpha)=\sum_iq_i\alpha_i-3(b\alpha_1\alpha_2\alpha_3)^{1/3}$. Let
$g=(\alpha_1\alpha_2\alpha_3)^{1/3}$, $w=1/\alpha$ and $s=1+2b$. The Hessian of $g$ in all 3
coordinates is $-gM$ with $M=\frac13\operatorname{diag}(w)^2-\frac19ww^\top$, so
$\nabla^2I=3b^{1/3}g\,J^\top MJ$, where $J$ is the Jacobian of
$(\alpha_2,\alpha_3)\mapsto(1-\alpha_2-\alpha_3,\alpha_2,\alpha_3)$. At $\alpha^*$,
$3b^{1/3}g=3b/s$ and $w=s(1,y,y)$ with $y=1/b$. So $J^\top MJ$ has diagonal entries
$\frac29s^2(1+y+y^2)$ and off-diagonal entries $\frac19s^2(2+2y-y^2)$, and
$\det J^\top MJ=\frac{s^4}{81}\cdot3y^2(2+y)^2=\frac{s^6}{27b^4}$. Hence
$\det\nabla^2I=s^4/(3b^2)$, and Proposition~\ref{prop:constant}(a) gives $K_{\{1,2,3\}}$. The
exact parts of (i) are checked in Lean.

(ii) Let $D_1=\log M_{\{1,2\}}-\log M_{\{1\}}=\log(bh)+(N-2)\log(1-h)-(N-1)\log(1-bh)$. Then
$D_1'(h)=\frac1h-\frac{N-2}{1-h}+\frac{b(N-1)}{1-bh}>0$ for $bh<1$, while $D_1\to-\infty$ as
$h\to0$ and $D_1\to\infty$ as $bh\to1$. So there is exactly one root (this is checked in
Lean). Expanding, $D_1=\log b+\log T-L+(b-1)T+O(T^2/N)$, which is positive at $T=L/(b-1)$ and
negative at $T=L/(2(b-1))$. So the root lies between them, and $\log T=\log L-\log(b-1)+O(\log
L/L)$ gives the expansion.

(iii) Here $\log M_{\{1,2\}}=\log b+\log T-L-T+O(T^2/N)$, and Theorem~\ref{thm:facemax} for the
full face ($k=3$, $\lambda=0$, $\mu(F)=1$) gives $\log M_{\{1,2,3\}}=-2L+\log T+\log
K_{\{1,2,3\}}+O(1/T+T^2/N)$. So on $[L/2,2L]$
the difference is $L-T+\log b-\log K_{\{1,2,3\}}+O(1/L)$, which gives (iii).

(iv) Lemma~\ref{lem:upper} with $\theta=q$ gives $M_{\{1,2,3\}}\le C(T/N)^2$. For $T<T^{(1)}_N$,
$M_{\{1\}}>M_{\{1,2\}}$ by (ii), and $M_{\{1\}}\ge N^{-b/(b-1)-o(1)}$ since $T^{(1)}_N\le
L/(b-1)$. This beats $C(T/N)^2$ because $b/(b-1)<2$. For $T>T^{(1)}_N$,
$M_{\{1,2\}}>M_{\{1\}}$. For $T\le L/2$, $\log M_{\{1,2\}}-\log M_{\{1,2,3\}}\ge L-T-\log T-C>0$.
On $[L/2,2L]$ the sign is that of the formula in (iii), which is positive at $L/2$ and negative at
$2L$. For $2L\le T\le N^{1/3}$ the same formula is at most $-L+C<0$.
\end{proof}

\begin{remark}[the next term, proof omitted]\label{rem:nextterm}
Carrying the Laplace argument behind Theorem~\ref{thm:local} one term further gives the factor
$1+c(\alpha)/T+O(T^{-3/2})$ there. Here $c(\alpha)$ is the coefficient of $1/T$ in the next step
of Lemma~\ref{lem:laplace} for the integral~\eqref{eq:A-fT}, divided by $\chi(\alpha,0)$. It is a
Gaussian moment of the $\varphi$-derivatives of $\Psi$ of orders $3$ and $4$ and of $\chi$ of
orders $1$ and $2$ at $\varphi=0$. Then~\eqref{eq:logmax} gains the term $c_F/T$, with
$c_F=c(\alpha^*_F)+\frac12\gamma^\top\Sigma_{\alpha^*_F}\gamma$ and
$\gamma=\nabla\log C_F(\alpha^*_F)$, and the error becomes $O(T^{-3/2}+T^2/N)$. So at every
crossing
\[
T_N\Bigl(\Delta L-\frac\Delta2\log T_N-\delta T_N-\log\frac{K_F}{K_G}\Bigr)
=c_F-c_G+O\bigl(T_N^{-1/2}+T_N^3/N\bigr).
\]
Here $c_{\{i\}}=0$. Comparing with the terms of order $1/z$ in the expansions of~\cite{paperA}
and~\cite{paperB} gives $c_F=-1/(8q_{ij})$ for a pair with $q_{ij}=q_{ji}$, $q_i=q_j$ and
$\mu_i=\mu_j$, and $c_F=-(k-1)/8$ for a face of $k$ states when $Q=J-dI$ and $\mu$ is constant on
$F$. We omit the details.
\end{remark}
